% =====================================================================
% supplementary.tex
% Heavy derivations and extension analyses supporting the main text:
%   - the correlated no-false-alarm integral and its monotonicity
%   - the closed-form value-directed-attention escape threshold
%   - symmetric recovery at unit benefit/cost ratio
%   - the power-mean conservation family
% =====================================================================

\appendix
\section{Supplementary derivations and extensions}
\label{sec:appendix}

This material derives the results used in the main text and develops the
model's scope. Section~\ref{sec:supp-decorrelation} derives the
one-dimensional integral for the no-false-alarm probability under
correlated decision noise and establishes its monotonicity in the
correlation. Section~\ref{sec:supp-escape} derives the closed-form
escape threshold for value-directed attention and extends it across the
correlation family. Section~\ref{sec:supp-symmetric} records the exact
recovery of the symmetric model at unit benefit/cost ratio.
Section~\ref{sec:supp-conservation} develops the one-parameter
conservation family that generalises the benefit/cost weighting and
establishes the invariances on which the headline results rest.
Throughout, $\varphi$ is the standard normal density, $\Phinorm$ its
cumulative distribution function, and $b_i := c_i + \dprime_i/2$ the
$z$-score the no-change decision variable must stay below at location
$i$.

% ---------------------------------------------------------------------
\subsection{The correlated no-false-alarm probability and its
monotonicity}
\label{sec:supp-decorrelation}

The only cross-location term in the expected reward
\eqref{eq:expected-reward} is the probability $\PnoFA$ that no location
produces a false alarm on a no-change trial
(Section~\ref{sec:model-reward}). Under independent per-location
decision variables this is the product
\eqref{eq:pnofa-indep}; here we derive its exact form when the decision
variables share an equicorrelation $\corr$ and show that admitting
positive correlation can only raise it.

\paragraph{The exact one-dimensional reduction.}
Let the no-change decision variables be jointly Gaussian with
standardised marginals and equicorrelation, as in
\eqref{eq:equicorr-cov}. The one-factor representation
\begin{equation}
  X_i \;=\; \mu_i \,+\, \sqrt{\corr}\;Z \,+\, \sqrt{1 - \corr}\;\varepsilon_i,
  \qquad
  Z, \varepsilon_1, \dots, \varepsilon_\Nloc
  \stackrel{\text{iid}}{\sim} \mathcal{N}(0, 1),
  \label{eq:one-factor}
\end{equation}
with $\mu_i = -\dprime_i/2$, reproduces $\operatorname{Var}(X_i) = 1$
and $\operatorname{Cov}(X_i, X_j) = \corr$ for $i \neq j$, and makes the
$X_i$ conditionally independent given the shared latent $Z$. Conditional
on $Z = z$, the per-location no-false-alarm event $\{X_i \le c_i\}$ has
probability
\begin{equation}
  \Pr(X_i \le c_i \mid Z = z)
  \;=\;
  \Phinorm\!\left(\frac{c_i - \mu_i - \sqrt{\corr}\,z}{\sqrt{1 - \corr}}\right)
  \;=\;
  \Phinorm\!\left(\frac{b_i - \sqrt{\corr}\,z}{\sqrt{1 - \corr}}\right).
  \label{eq:conditional-noFA}
\end{equation}
Multiplying across the conditionally independent locations and
integrating $Z$ out against $\varphi$ gives the exact one-dimensional
integral \eqref{eq:pnofa-rho}. Equation~\eqref{eq:pnofa-rho} holds for
any $\Nloc$ and any $\corr \in [0, 1)$: equicorrelation collapses the
$\Nloc$-dimensional orthant integral to a single dimension, with no
multivariate-normal cumulative distribution function required. At
$\corr = 0$ the integrand of \eqref{eq:pnofa-rho} is constant in $z$ and
the integral reduces to the independent product
\eqref{eq:rho-zero-recovery}. The model evaluates
\eqref{eq:pnofa-rho} by $64$-node Gauss--Hermite quadrature; with the
substitution $z = \sqrt{2}\,t$ the integrand carries weight $e^{-t^2}$,
and the $64$-node rule agrees with a $128$-node reference to better than
$10^{-15}$ across the operating band $\corr \in
\{0, 0.05, 0.1, 0.2, 0.3, 0.4\}$.

\paragraph{Monotonicity in the correlation.}
\begin{proposition}[Orthant monotonicity]
\label{prop:orthant-monotone}
For jointly Gaussian decision variables with standardised marginals and
equicorrelation $\corr \in [0, 1)$, the no-false-alarm probability
$\PnoFA(\corr)$ of \eqref{eq:pnofa-rho} is non-decreasing in $\corr$ at
every fixed threshold pair $(b_c, b_u)$, with strict increase for
$\corr > 0$ at finite thresholds:
\begin{equation}
  \PnoFA(\corr) \;\ge\; \PnoFA(0)
  \;=\; \Phinorm(b_c)\,\Phinorm(b_u)^{\Nloc - 1}
  \qquad \forall\,\corr \in [0, 1).
  \label{eq:orthant-monotone}
\end{equation}
\end{proposition}

\begin{proof}
Equation~\eqref{eq:orthant-monotone} is the equicorrelated case of
Slepian's monotonicity inequality for Gaussian orthant probabilities,
which states that $\Pr\!\big(\bigcap_i \{X_i \le c_i\}\big)$ is
non-decreasing in each off-diagonal correlation entry at fixed
thresholds \cite{Slepian1962}; applying it to every pair correlation
simultaneously gives monotonicity in the common value $\corr$. A modern
treatment is given in \cite{Tong1990}.
\end{proof}

Positive correlation relaxes the multiple-comparisons pressure that, in
the independent model, drives the $\Nloc$-fold no-false-alarm product
toward zero whenever any single $\Phinorm(b_i)$ falls below one. The
independent corner $\corr = 0$ is therefore the stiffest-penalty
boundary of the family: it maximises the aggregate cost of holding
$\Nloc$ false-alarm rates below threshold at once. Because the
change-trial bracket of \eqref{eq:expected-reward} involves only
marginal hit rates and is independent of $\corr$, and the no-change
bracket scales linearly in $\PnoFA(\corr)$ with the non-negative
coefficient $\tfrac{1}{2}\CR$, Proposition~\ref{prop:orthant-monotone}
gives a pointwise reward bound that survives the supremum over criteria.

\begin{corollary}[Per-policy reward monotonicity]
\label{cor:policy-monotone}
At any fixed allocation $\alpha$, the criterion-optimised expected reward
$\Rew^\star(\alpha; \corr) := \sup_{c_c, c_u}
\mathbb{E}[\Rew](\alpha, c_c, c_u; \corr)$ is non-decreasing in $\corr$.
Consequently each of the nested policies~P1--P4 of
Section~\ref{sec:model-policies} has its optimal reward non-decreasing
in $\corr$.
\end{corollary}

Corollary~\ref{cor:policy-monotone} bounds the reward \emph{levels} but
not their \emph{difference}: the value-directed-attention benefit
$\VDA(\corr) = \Rpone(\corr) - \Rptwo(\corr)$ is the gap between two
quantities that both rise with $\corr$, and such a gap can move in
either direction. The sign of $\partial \VDA / \partial \corr$ is
therefore a property of the model to be determined, not a corollary of
monotonicity. Two mechanisms compete. Raising $\corr$ inflates
$\PnoFA(\corr)$ most at liberal criteria, flattening the no-change
bracket across $(c_c, c_u)$ and \emph{devaluing} the criterion lever;
this is the channel that drives the criterion fraction $\CF(\corr)$
downward (Section~\ref{sec:results-criterion}). At the same time, when
the optimal allocation concentrates attention on the cued location the
uncued sensitivity falls and the uncued false-alarm rate rises,
compounded across $\Nloc - 1$ locations through the no-false-alarm
product; Proposition~\ref{prop:orthant-monotone} shows this aggregate
penalty is \emph{relaxed} under correlation, making concentration
cheaper and \emph{amplifying} $\VDA$. The criterion-devaluation channel
dominates in the cost-dominant regime (small $\Rsens$) and the
concentration-relaxation channel dominates in the benefit-dominant
regime (large $\Rsens$), so the sign of $\partial \VDA / \partial \corr$
flips with $\Rsens$ — the model property reported across the parameter
plane in Sections~\ref{sec:results-vda-nonmonotonic}
and~\ref{sec:results-graded}.

% ---------------------------------------------------------------------
\subsection{Closed-form escape threshold for value-directed attention}
\label{sec:supp-escape}

Here we derive the escape threshold $\rdagger(\val)$ of
\eqref{eq:r-dagger}--\eqref{eq:K-u} and extend it across the correlation
family. The threshold is the value of the benefit/cost ratio below which
the value-aware optimum is still pinned at uniform allocation
$\alpha = 1/\Nloc$, and above which the optimiser begins to move
sensitivity toward the cued location.

\paragraph{Setup at the boundary.}
At the uniform allocation $\alpha = 1/\Nloc$ the asymmetric sensitivity
rule of Section~\ref{sec:model-mapping} collapses:
$\dprime_c = \dprime_u = \dprime_{\mathrm{base}} := \dprimemax\,
f(1/\Nloc)$, independent of $(\Rsens, \corr)$, because the
gain bracket $\dprimemax\, f(1/\Nloc) - \dprime_{\mathrm{base}}$ vanishes
before being weighted by $\benefit(\Rsens)$ or $\cost(\Rsens)$. The
criterion optimum at this allocation is therefore the joint maximiser of
the expected reward over a fixed-$\dprime$ landscape that depends on
$\corr$ only through the correlated no-false-alarm term:
\begin{equation}
  \bigl(c_c^\star(\corr),\, c_u^\star(\corr)\bigr)
  \;:=\;
  \operatorname*{arg\,max}_{(c_c, c_u)}\,
  \Rew\bigl(\dprime_{\mathrm{base}},\, \dprime_{\mathrm{base}},\,
           c_c,\, c_u;\, \valid, \val, \Nloc, \CR, \corr\bigr).
  \label{eq:boundary-crit}
\end{equation}
At the headline cell $(\Nloc, \valid, \val) = (4, 0.5, 5)$ with
$\CR = \valid\val + (1 - \valid) = 3.0$, the optimal pair
$(c_c^\star, c_u^\star) = (0.10, 1.75)$ at $\corr = 0$ shifts to
$(0.05, 1.80)$ at $\corr = 0.2$: the cued criterion becomes more liberal
(a cued false alarm is partially explained away by correlated false
alarms at the other locations) and the uncued criterion more
conservative (an uncued false alarm now also degrades the cued
no-false-alarm joint event).

\paragraph{The correlation-aware gradient integrals.}
Differentiating $\PnoFA(\corr)$ of \eqref{eq:pnofa-rho} under the
integral sign — the integrand and its $\dprime$-derivatives are bounded,
so dominated convergence applies — gives two one-dimensional
Gauss--Hermite integrals,
\begin{align}
  I_c(b_c, b_u; \corr, \Nloc)
  &\;:=\;
  \int_{-\infty}^{\infty}
  \frac{1}{\sqrt{1 - \corr}}\,
  \varphi\!\left(\frac{b_c - \sqrt{\corr}\,z}{\sqrt{1 - \corr}}\right)
  \Phinorm\!\left(\frac{b_u - \sqrt{\corr}\,z}{\sqrt{1 - \corr}}\right)^{\Nloc - 1}
  \varphi(z)\,dz,
  \label{eq:gh-grad-c}
  \\[2pt]
  I_u(b_c, b_u; \corr, \Nloc)
  &\;:=\;
  \int_{-\infty}^{\infty}
  \frac{1}{\sqrt{1 - \corr}}\,
  \Phinorm\!\left(\frac{b_c - \sqrt{\corr}\,z}{\sqrt{1 - \corr}}\right)
  \Phinorm\!\left(\frac{b_u - \sqrt{\corr}\,z}{\sqrt{1 - \corr}}\right)^{\Nloc - 2}
  \varphi\!\left(\frac{b_u - \sqrt{\corr}\,z}{\sqrt{1 - \corr}}\right)
  \varphi(z)\,dz.
  \label{eq:gh-grad-u}
\end{align}
At $\corr = 0$ these reduce to the independent products
\begin{equation*}
  I_c^0 = \varphi(b_c)\,\Phinorm(b_u)^{\Nloc - 1},
  \qquad
  I_u^0 = \Phinorm(b_c)\,\Phinorm(b_u)^{\Nloc - 2}\,\varphi(b_u),
\end{equation*}
and they share the same Gauss--Hermite nodes used to evaluate
$\PnoFA(\corr)$ itself.

\paragraph{Boundary first-order condition.}
The chain-rule expansion of $\partial \Rew / \partial\alpha$ at
$\alpha = 1/\Nloc^{+}$ collects the asymmetric criterion booking into
\begin{equation}
  \frac{\partial \Rew}{\partial\alpha}\bigg|_{1/\Nloc^{+}}\!(\corr)
  \;=\;
  \dprimemax\, f'(1/\Nloc)\,
  \Bigl[\,\benefit(\Rsens)\,K_c(\val; \corr)
        \;-\; \cost(\Rsens)\,\frac{K_u(\val; \corr)}{\Nloc - 1}\,\Bigr],
  \label{eq:boundary-foc-rho}
\end{equation}
with correlation-aware coefficients
\begin{align}
  K_c(\val; \corr)
  &= \tfrac{1}{4}\Big[\valid\,\val\,
       \varphi\!\big(\tfrac{\dprime_b}{2} - c_c^\star(\corr)\big)
       \;+\; \CR\,
       I_c\bigl(b_c^\star(\corr), b_u^\star(\corr); \corr, \Nloc\bigr)\Big],
  \label{eq:Kc-rho}
  \\[2pt]
  K_u(\val; \corr)
  &= \tfrac{1}{4}\Big[(1 - \valid)\,
       \varphi\!\big(\tfrac{\dprime_b}{2} - c_u^\star(\corr)\big)
       \;+\; (\Nloc - 1)\,\CR\,
       I_u\bigl(b_c^\star(\corr), b_u^\star(\corr); \corr, \Nloc\bigr)\Big],
  \label{eq:Ku-rho}
\end{align}
where $b_i^\star(\corr) := c_i^\star(\corr) + \dprime_{\mathrm{base}}/2$,
$\dprime_b \equiv \dprime_{\mathrm{base}}$, and $f'(a) = (1 - f_0)\,h'(a)$.
The change-trial density terms depend on $\corr$ only through
$c_i^\star(\corr)$; the no-false-alarm terms $I_c, I_u$ carry the full
correlation dependence of the quadrature. Both $K_c$ and $K_u$ are
positive.

\begin{proposition}[Escape threshold under equicorrelation]
\label{prop:escape-rho}
Fix $(\Nloc, \valid, \val, \dprimemax, f_0, h, \CR, \corr)$ with
$\corr \in [0, 1)$, and let $(c_c^\star(\corr), c_u^\star(\corr))$ solve
\eqref{eq:boundary-crit}. Then the value-aware optimum satisfies
$\alphacued^\star = 1/\Nloc$ for $\Rsens \le \rdagger(\val; \corr)$ and
$\alphacued^\star > 1/\Nloc$ for $\Rsens > \rdagger(\val; \corr)$, with
\begin{equation}
  \boxed{\;
    \rdagger(\val; \corr)
    \;=\;
    \frac{K_u(\val; \corr)}{(\Nloc - 1)\,K_c(\val; \corr)}.
  \;}
  \label{eq:r-dagger-rho}
\end{equation}
\end{proposition}

\begin{proof}
Setting \eqref{eq:boundary-foc-rho} to zero and dividing out the
strictly positive prefactor
$\dprimemax\, f'(1/\Nloc)\,\cost(\Rsens) > 0$ gives
$\Rsens\, K_c(\val; \corr) = K_u(\val; \corr)/(\Nloc - 1)$, using
$\benefit(\Rsens)/\cost(\Rsens) = \Rsens$ (which holds for every member
of the conservation family of Section~\ref{sec:supp-conservation});
solving for $\Rsens$ gives \eqref{eq:r-dagger-rho}. The corner maximum at
$\alpha = 1/\Nloc$ becomes a saddle as $\Rsens$ crosses the threshold; a
one-sided finite-difference check at
$\Rsens = \rdagger(\val; \corr) \pm 0.05$ confirms the boundary slope
changes sign at every $(\val, \corr) \in \{1, 2, 3\} \times \{0, 0.2\}$
tested.
\end{proof}

\paragraph{Recovery at zero correlation.}
At $\corr = 0$ the integrand of \eqref{eq:gh-grad-c} loses its
$z$-dependence and $I_c \to I_c^0$, $I_u \to I_u^0$; the criterion
optimum \eqref{eq:boundary-crit} reduces to the corresponding optimum at
$\corr = 0$, so $K_c(\val; 0) = K_c(\val)$ and $K_u(\val; 0) = K_u(\val)$
of \eqref{eq:K-c}--\eqref{eq:K-u}, and
Proposition~\ref{prop:escape-rho} reduces to \eqref{eq:r-dagger}. The
recovery is structural — the integrals collapse term by term — and the
threshold computed from \eqref{eq:r-dagger-rho} at $\corr = 0$
reproduces Table~\ref{tab:r-dagger-family} across $\val \in \{1, 2, 3,
5, 8, 10\}$ to within the criterion-grid resolution
($\le 5 \times 10^{-4}$ in $\rdagger$).

\paragraph{Drift under correlation.}
Evaluating \eqref{eq:r-dagger-rho} at the headline cell quantifies how
the escape threshold moves when correlation is admitted
(Table~\ref{tab:r-dagger-rho-drift}).
%
\begin{table}[t]
  \centering
  \small
  \begin{tabular}{r r r r r c}
    \toprule
    $\val$ &
    $\rdagger(\val; 0)$ &
    $\rdagger(\val; 0.2)$ &
    $\Delta\rdagger$ &
    $\%\,\Delta\rdagger$ &
    sign-match $r^\star$? \\
    \midrule
    $1$  & $0.343$ & $0.354$ & $+0.0103$ & $+3.0\%$  & --- (reference) \\
    $2$  & $0.168$ & $0.173$ & $+0.0051$ & $+3.0\%$  & \checkmark \\
    $3$  & $0.099$ & $0.113$ & $+0.0132$ & $+13.2\%$ & \checkmark \\
    $5$  & $0.050$ & $0.058$ & $+0.0077$ & $+15.2\%$ & \checkmark \\
    $8$  & $0.022$ & $0.029$ & $+0.0066$ & $+29.7\%$ & \checkmark \\
    $10$ & $0.016$ & $0.019$ & $+0.0030$ & $+18.7\%$ & \checkmark \\
    \bottomrule
  \end{tabular}
  \caption{Drift of the escape threshold under correlation at the
  headline cell $(\Nloc, \valid, \dprimemax, f_0, h) =
  (4,\, 0.5,\, 2,\, 0.5,\, \sqrt{\;\cdot\;})$, value-coupled correct
  rejection, additive conservation. $\rdagger(\val; \corr)$ from
  \eqref{eq:r-dagger-rho}; $\Delta\rdagger := \rdagger(\val; 0.2) -
  \rdagger(\val; 0)$. The drift is positive at every cued value,
  including the value-blind reference $\val = 1$: the regime in which the
  optimum is pinned at uniform allocation \emph{widens} as correlation
  grows. The sign of $\Delta\rdagger$ matches the sign of the empirical
  peak drift $\Delta r^\star$ of Table~\ref{tab:rho-sensitivity} at all
  five $\val \ne 1$ rows.}
  \label{tab:r-dagger-rho-drift}
\end{table}
%
Two features follow. First, $\Delta\rdagger(\val) > 0$ at every $\val$:
the lower edge of the active band moves up under correlation across the
whole value family, and because the peak of the value-directed-attention
benefit cannot lie below the escape threshold, the peak must drift up
with it — the behaviour reported in
Section~\ref{sec:results-vda-nonmonotonic}. Second, the absolute drift in
the threshold is bounded by $\le 0.014$ across the family, while the
empirical peak drift in Table~\ref{tab:rho-sensitivity} reaches $0.13$ at
$\val = 2$: the peak responds not only to the lower edge but also to the
upward drift of the value-blind reference threshold (which widens the
band's upper edge) and to how steeply the optimal allocation climbs above
uniform once it escapes. Equation~\eqref{eq:r-dagger-rho} pins the
lower-edge drift exactly; a closed form for the peak location itself,
which would additionally model the climb of $\alphacued^\star$ above the
threshold, is a natural further step. The two gradient integrals
\eqref{eq:gh-grad-c}--\eqref{eq:gh-grad-u} make the concentration-cost
channel of Section~\ref{sec:supp-decorrelation} analytic on the lower
edge: the same one-factor reduction that renders $\PnoFA(\corr)$
tractable renders the boundary slope tractable.

% ---------------------------------------------------------------------
\subsection{Symmetric recovery at unit benefit/cost ratio}
\label{sec:supp-symmetric}

At $\Rsens = 1$ the asymmetric benefit/cost weighting
$\benefit(\Rsens) = 2\Rsens/(\Rsens + 1)$,
$\cost(\Rsens) = 2/(\Rsens + 1)$ collapses to the symmetric special case
$\benefit = \cost = 1$, in which a single shared transfer function drives
both the cued and the uncued sensitivity. We record the recovery at the
strength the model supports: a real-number identity that holds
universally, and a bit-exact identity that holds at the validation
configuration.

\paragraph{The collapse as a real-number identity.}
The asymmetric sensitivity rule of Section~\ref{sec:model-mapping} writes
$\dprime_i(\alpha) = \dprime_{\mathrm{base}} + s_i(\Rsens)\,
(\dprimemax\, f(\alpha) - \dprime_{\mathrm{base}})$ with
$s_i \in \{\benefit, \cost\}$ according to whether location $i$ is
attended; at $\Rsens = 1$ both slopes are unit and the expression
collapses to $\dprime_i(\alpha) = \dprimemax\, f(\alpha)$, the symmetric
rule. Because the downstream signal-detection, expected-reward, and
criterion-optimisation stack is a deterministic function of the
$\dprime$ arrays alone, identical arrays force identical optima.

\begin{proposition}[Real-number recovery at unit ratio]
\label{prop:symmetric-recovery}
For every $(\valid, \val, \Nloc, \dprimemax, f_0, h)$ in the model's
domain and at any correlation $\corr \in [0, 1)$, the asymmetric model
evaluated at $\Rsens = 1$ produces the same $\dprime$ arrays as the
symmetric special case $\benefit = \cost = 1$. Consequently the optimal
cued share, the optimal cued and uncued criteria, the optimal reward, and
the criterion fraction all match the symmetric baseline as real-number
identities.
\end{proposition}

This is a property of the algebra rather than of the numerical
implementation: the cancellation
$\dprime_{\mathrm{base}} - \dprime_{\mathrm{base}} = 0$ is the entire
content of the recovery.

\paragraph{Bit-exact identity at the validation configuration.}
A sharper statement holds at the configuration used to validate the
model. The asymmetric code path at $\Rsens = 1$ forms
$\mathrm{fl}\big(a + \mathrm{fl}(1.0 \cdot \mathrm{fl}(x - a))\big)$ with
$a = \dprime_{\mathrm{base}}$ and $x = \dprimemax\, f(\alpha)$. Three
facts compose: (i) $\benefit(1)$ and $\cost(1)$ are the exact binary
float $1.0$, so the multiplicative slope round-trips identically; (ii)
Sterbenz's lemma~\cite{Sterbenz1974, Goldberg1991} guarantees
$\mathrm{fl}(x - a) = x - a$ exactly whenever $a/2 \le x \le 2a$; and
(iii) at the validation configuration
$(\Nloc, \dprimemax, f_0, h) = (4, 2.0, 0.5, \sqrt{\;\cdot\;})$, value-coupled
correct rejection, the baseline is
$a = \dprimemax(f_0 + (1 - f_0)\,h(1/\Nloc)) = 1.5$ and the swept range
$x = \dprimemax\, f(\alpha) \in [1.0, 2.0]$ lies inside the Sterbenz band
$[a/2, 2a] = [0.75, 3.0]$ at every grid point.

\begin{proposition}[Bit-exact symmetric recovery]
\label{prop:bitexact-recovery}
Under value-coupled correct rejection at the validation configuration,
the asymmetric implementation evaluated at $\Rsens = 1$ returns
$\dprime$ arrays bit-identical to the symmetric implementation; hence the
optimal cued share and optimal reward match it to zero difference, not
merely to machine epsilon.
\end{proposition}

The hypothesis of Proposition~\ref{prop:bitexact-recovery} is sufficient
but not necessary, and fails once the smallest swept output $\dprimemax\,
f_0$ drops below $\dprime_{\mathrm{base}}/2$, i.e.\ once
\begin{equation}
  f_0 \;<\; \frac{h(1/\Nloc)}{1 + h(1/\Nloc)} \;=\; \frac{1}{3}
  \qquad (h = \sqrt{\;\cdot\;},\; \Nloc = 4).
  \label{eq:sterbenz-threshold}
\end{equation}
Off-band, bit-identity is lost and the two paths drift at the order of a
unit in the last place. The universal statement is therefore that the
symmetric recovery is exact to machine precision; the bit-for-bit
identity is a structural guarantee of the chosen configuration rather
than a property of the model.

\paragraph{Unit ratio is the smooth centre of the family.}
The collapse at $\Rsens = 1$ is the interior of a smooth, two-sided
family, not a removable singularity. The slope ratio
$\benefit(\Rsens)/\cost(\Rsens) = \Rsens$ is continuous and
differentiable through $\Rsens = 1$, with
$d\benefit/d\Rsens|_{1} = 1/2$ and $d\cost/d\Rsens|_{1} = -1/2$. A
continuity check at $\Rsens = 1 \pm \epsilon$ confirms that the maximum
reward difference against the symmetric model scales linearly in
$|\Rsens - 1|$, with slope $\approx 0.084$ reward units per unit shift.
Unit ratio is the balanced point of the enhancement-versus-suppression
trade-off (Definition~\ref{def:three-levers}), the natural null about
which the value-directed-attention effects of
Section~\ref{sec:results} are organised. The recovery is invariant to the
conservation rule: at $\Rsens = 1$ the ratio identity
$\benefit/\cost = 1$ forces $\benefit = \cost$, and any conservation
constraint in the family of Section~\ref{sec:supp-conservation} then
forces $\benefit = \cost = 1$, so
Propositions~\ref{prop:symmetric-recovery}--\ref{prop:bitexact-recovery}
hold under every conservation order without modification
(Corollary~\ref{cor:symmetric-invariance}).

% ---------------------------------------------------------------------
\subsection{The power-mean conservation family}
\label{sec:supp-conservation}

The benefit/cost weights satisfy the asymmetry ratio
$\benefit/\cost = \Rsens$ together with a conservation constraint; the
main text uses the additive rule $\benefit + \cost = 2$
(Eq.~\ref{eq:beta-gamma}). This subsection embeds that rule in a
one-parameter family, derives the family's monotonicity, and establishes
the two invariances on which the escape threshold and the symmetric
recovery rest.

\paragraph{The family and its closed-form weights.}
For $p \in \mathbb{R}$, the power (Hölder) mean of order $p$ on positive
reals $(x, y)$ is
\begin{equation}
  M_p(x, y) \;=\;
  \begin{cases}
    \displaystyle \left( \frac{x^p + y^p}{2} \right)^{1/p}, & p \ne 0, \\[4pt]
    \displaystyle \sqrt{x\,y}, & p = 0,
  \end{cases}
  \label{eq:power-mean}
\end{equation}
with limiting cases $M_{-\infty} = \min$ and $M_{+\infty} = \max$
\cite{HLP1934}. The conservation family is the constraint
$M_p(\benefit, \cost) = 1$ together with $\benefit/\cost = \Rsens$.
Substituting $\benefit = \Rsens\,\cost$ for $p \ne 0$ gives the
closed-form weights
\begin{equation}
  \boxed{\;
    \cost(\Rsens; p) \;=\;
    \left( \frac{2}{\Rsens^{\,p} + 1} \right)^{1/p},
    \qquad
    \benefit(\Rsens; p) \;=\; \Rsens\,\cost(\Rsens; p),
  \;}
  \label{eq:power-mean-weights}
\end{equation}
with the geometric-mean limit $\cost(\Rsens; 0) = 1/\sqrt{\Rsens}$,
$\benefit(\Rsens; 0) = \sqrt{\Rsens}$. The case $p = 1$ recovers the
additive rule $\benefit = 2\Rsens/(\Rsens + 1)$,
$\cost = 2/(\Rsens + 1)$; $p = 0$ is the multiplicative
(geometric-mean) endpoint $\benefit\cost = 1$. Three identities hold at
every $p$: $\benefit/\cost = \Rsens$ by construction;
$\benefit(1; p) = \cost(1; p) = 1$; and
$\operatorname{sgn}(\benefit - \cost) = \operatorname{sgn}(\Rsens - 1)$.

\paragraph{Monotonicity in the conservation order.}
The classical power-mean inequality makes $p \mapsto M_p(x, y)$ strictly
increasing for fixed $x \ne y$ \cite{HLP1934}. Under the constrained form
$M_p(\benefit, \cost) = 1$, this translates into an exact identity on
$\partial \ln \cost / \partial p$. Writing
$L(p) = \ln \cost(\Rsens; p) = (1/p)\ln\!\big(2/(\Rsens^p + 1)\big)$ and
$\theta_p := \Rsens^p / (\Rsens^p + 1)$, direct differentiation gives
\begin{equation}
  \boxed{\;
    \frac{\partial}{\partial p}\,\ln \cost(\Rsens; p)
    \;=\;
    -\frac{1}{p^2}\,
    \DKL\!\bigl(\,\mathrm{Bern}(\theta_p)\,\big\|\,
                    \mathrm{Bern}(1/2)\,\bigr),
  \;}
  \label{eq:hlp-kl}
\end{equation}
where the bracket of the expansion,
$\ln(2/(\Rsens^p + 1)) + p\,\theta_p \ln \Rsens$, equals
$\theta_p \ln(2\theta_p) + (1 - \theta_p)\ln(2(1 - \theta_p)) =
\DKL(\mathrm{Bern}(\theta_p)\,\|\,\mathrm{Bern}(1/2))$. Since
$\DKL \ge 0$ with equality only at $\theta_p = 1/2$ (i.e.\ $\Rsens = 1$),
\eqref{eq:hlp-kl} yields, for every $\Rsens \ne 1$ and every
$p \ne 0$,
\begin{equation}
  \frac{\partial \cost}{\partial p}(\Rsens; p) \;<\; 0,
  \qquad
  \frac{\partial \benefit}{\partial p}(\Rsens; p)
  \;=\; \Rsens\,\frac{\partial \cost}{\partial p} \;<\; 0.
  \label{eq:mono-sign}
\end{equation}
Continuity through $p = 0$ holds by series expansion, with
$\lim_{p \to 0}\partial \ln \cost / \partial p = -(\ln \Rsens)^2 / 8$.
Both weights therefore increase strictly as the conservation order
decreases, at every $\Rsens \ne 1$. A finite-difference cross-check of
\eqref{eq:hlp-kl} across a range of $(\Rsens, p)$ pairs agrees to
$\le 1.5 \times 10^{-10}$.

\paragraph{The symmetric corner and the recovery at unit ratio.}
\begin{proposition}[Symmetric-corner identity]
\label{prop:symmetric-corner}
For every $p \in \mathbb{R}$ (including $p = 0$),
$\benefit(1; p) = \cost(1; p) = 1$.
\end{proposition}

\begin{proof}
At $\Rsens = 1$, $\Rsens^p = 1$ for every $p$, so
\eqref{eq:power-mean-weights} gives $\cost = (2/2)^{1/p} = 1$ and
$\benefit = 1$ for $p \ne 0$; the geometric-mean limit gives
$\cost = 1/\sqrt{1} = 1$, $\benefit = 1$ at $p = 0$.
\end{proof}

\begin{corollary}[Conservation-form invariance of the symmetric recovery]
\label{cor:symmetric-invariance}
The symmetric recovery of Proposition~\ref{prop:symmetric-recovery} holds
at every conservation order: at $\Rsens = 1$ and any $p$, the asymmetric
$\dprime$-map collapses pointwise to the symmetric baseline,
$\dprime_c(\alpha; 1, p) = \dprime_u\big((1 - \alpha)/(\Nloc - 1); 1, p\big)
= \dprimemax\, f(\alpha)$.
\end{corollary}

\begin{proof}
By Proposition~\ref{prop:symmetric-corner},
$\benefit(1; p) = \cost(1; p) = 1$ at every $p$, so the $\dprime$-map at
$\Rsens = 1$ reads $\dprime_c(\alpha; 1, p) = \dprime_{\mathrm{base}} +
1 \cdot (\dprimemax\, f(\alpha) - \dprime_{\mathrm{base}}) = \dprimemax\,
f(\alpha)$ with no $p$-dependence; the uncued counterpart is identical.
The optima then coincide with the symmetric baseline as in
Proposition~\ref{prop:symmetric-recovery}.
\end{proof}

\paragraph{Conservation-form invariance of the escape threshold.}
\begin{proposition}[Escape threshold is conservation-form invariant]
\label{prop:escape-invariance}
The escape threshold $\rdagger(\val) = K_u(\val)/[(\Nloc - 1) K_c(\val)]$
of \eqref{eq:r-dagger} is independent of the conservation order $p$.
\end{proposition}

\begin{proof}
\textbf{(Step 1)} At the uniform allocation $\alpha = 1/\Nloc$ the
$\dprime$-map reads
$\dprime_c(1/\Nloc; \Rsens, p) = \dprime_{\mathrm{base}} +
\benefit(\Rsens; p)\,(\dprimemax\, f(1/\Nloc) - \dprime_{\mathrm{base}})
= \dprime_{\mathrm{base}}$, because the bracket vanishes by the
definition $\dprime_{\mathrm{base}} := \dprimemax\, f(1/\Nloc)$. The
vanishing bracket annihilates any value of $\benefit(\Rsens; p)$, so
$\dprime_c$ at uniform allocation is $p$-independent; the uncued
per-location allocation is also $1/\Nloc$, so $\dprime_u =
\dprime_{\mathrm{base}}$ likewise. The pair
$(\dprime_{\mathrm{base}}, \dprime_{\mathrm{base}})$ is therefore
independent of $p$ (and of $\Rsens$).
\textbf{(Step 2)} At this fixed allocation the criterion-optimal pair
$(c_c^\star, c_u^\star)$ is the argmax of a reward functional whose only
$p$-dependence entered through $(\dprime_c, \dprime_u)$; by Step 1 those
are $p$-independent, so the criteria are too.
\textbf{(Step 3)} The gradient coefficients $K_c(\val), K_u(\val)$ are
the sensitivity partials of the reward evaluated at the $p$-independent
$(\dprime_{\mathrm{base}}, \dprime_{\mathrm{base}}, c_c^\star, c_u^\star)$,
and the functional form depends only on $(\valid, \val, \Nloc, \CR)$;
hence $K_c, K_u$ and their ratio $\rdagger(\val)$ are $p$-independent.
\end{proof}

The escape threshold is thus an algebraic invariant of the family: the
left edge of the active band is fixed across conservation orders, while
the peak inside the band — set at allocations $\alpha \ne 1/\Nloc$, where
the gain bracket no longer vanishes — varies with $p$ and is reported as
a band in Section~\ref{sec:results-vda-nonmonotonic}.

\paragraph{The criterion fraction under the conservation order.}
The dependence of the criterion fraction on $p$ admits a chain-rule sign
analysis. With $\CF = (\Rpthree - \Rpone)/(\Rpfour - \Rpone)$
(Eq.~\ref{eq:cf-def}) and the envelope theorem applied at each policy's
own optimum, the $p$-partials thread through the $\dprime$-map:
\begin{equation}
  \frac{\partial \Rpone}{\partial p}
  \;=\;
  \frac{\partial \Rpone}{\partial \dprime_c}\bigg|_{\star}
  \frac{\partial \dprime_c}{\partial p}
  \;+\;
  \frac{\partial \Rpone}{\partial \dprime_u}\bigg|_{\star}
  \frac{\partial \dprime_u}{\partial p},
  \label{eq:cf-chain-rule}
\end{equation}
with
\begin{equation}
  \frac{\partial \dprime_c}{\partial p}
  = \frac{\partial \benefit}{\partial p}\,
    \big(\dprimemax\, f(\alpha) - \dprime_{\mathrm{base}}\big),
  \qquad
  \frac{\partial \dprime_u}{\partial p}
  = \frac{\partial \cost}{\partial p}\,
    \big(\dprimemax\, f((1 - \alpha)/(\Nloc - 1)) - \dprime_{\mathrm{base}}\big).
  \label{eq:dprime-of-p}
\end{equation}
By \eqref{eq:mono-sign}, both $\partial \benefit/\partial p$ and
$\partial \cost/\partial p$ are negative for $\Rsens \ne 1$; for
$\alpha \ge 1/\Nloc$ the first bracket of \eqref{eq:dprime-of-p} is
non-negative and the second non-positive, so the cued and uncued
sensitivity channels move with opposite sign in $p$.

\begin{proposition}[Criterion-only reward is invariant at uniform
allocation]
\label{prop:p3-invariance}
At the criterion-only policy~P3, fixed at $\alpha = 1/\Nloc$,
$\partial \Rpthree / \partial p \equiv 0$.
\end{proposition}

\begin{proof}
By Step 1 of Proposition~\ref{prop:escape-invariance} the brackets in
\eqref{eq:dprime-of-p} vanish at $\alpha = 1/\Nloc$, so
$\partial \dprime_c/\partial p = \partial \dprime_u/\partial p = 0$ and
\eqref{eq:cf-chain-rule} gives $\partial \Rpthree/\partial p = 0$.
\end{proof}

Using Proposition~\ref{prop:p3-invariance}, the criterion-fraction
gradient simplifies to
\begin{equation*}
  \frac{\partial \CF}{\partial p}
  \;=\;
  \frac{-(\Rpfour - \Rpone)\,\partial_p \Rpone
        \;-\; (\Rpthree - \Rpone)\,(\partial_p \Rpfour - \partial_p \Rpone)}
       {(\Rpfour - \Rpone)^2},
\end{equation*}
whose sign reduces to a competition between $-\partial_p \Rpone \ge 0$
and $-(\partial_p \Rpfour - \partial_p \Rpone)$. The uniform inequality
$|\partial_p \Rpfour| \le |\partial_p \Rpone|$ would make
$\partial \CF/\partial p \le 0$ pointwise; this holds empirically across
the full parameter sweep, where swapping the additive rule ($p = 1$) for
the multiplicative endpoint ($p = 0$) gives $\Delta\CF \le 0$ in every
one of the $4{,}410$ valid cells, with no reverse flips, doubling the
fraction of cells with $\CF < 0.5$ from $4.0\%$ to $8.3\%$ while moving
the median by less than $0.005$. A closed-form proof of the uniform
inequality is left open: both $\partial_p \Rpone$ and
$\partial_p \Rpfour$ carry the moving optimum $\alpha^\star(p)$ that
the chain rule does not collapse.

\paragraph{Scope.}
The conservation-order analysis is carried at zero correlation; the
joint dependence on $(p, \corr)$ would compose the power-mean gradient
\eqref{eq:hlp-kl} with the correlation gradients
\eqref{eq:gh-grad-c}--\eqref{eq:gh-grad-u} and is a natural further step.
The analysis assumes a single global benefit/cost ratio and homogeneous
uncued allocation; relaxing either to per-location heterogeneity or to
the full $\Nloc$-dimensional allocation simplex leaves the
conservation weights $(\benefit, \cost)$ entering the per-location
$\dprime$-map unchanged, so the two relaxations compose with the
conservation order.
